% -*- mode: LaTeX; TeX-master: "lin-alg"; -*-
\oppgaver{1}

\begin{oppgave}
Tegn grafen til hver av disse likningene:
\begin{punkt}
$2x = 5y$
\end{punkt}
\begin{punkt}
$xy = 1$
\end{punkt}
\begin{punkt}
$\frac{x+y}{2} = \frac{7}{3}$
\end{punkt}
\begin{punkt}
$x^2 = y^2$
\end{punkt}
\begin{punkt}
$x^2 + y^2 + 2 = 0$
\end{punkt}
\medskip\noindent
Hvilke av dem er lineære?
\end{oppgave}

\begin{losning}
\begin{punkt}
Likningen $2x = 5y$ kan vi skrive om til
\[
2x - 5y = 0
\]
for å få den på den vanlige formen for lineære likninger (og dermed
ser vi med en gang at denne likningen er lineær).  For å gjøre det
enklere å tegne grafen, kan det være praktisk å «løse for~$y$», altså
å skrive om likningen til:
\[
y = \frac{2}{5} x
\]
Grafen blir altså en rett linje som går gjennom origo og har
stigningstall~$2/5$:
\begin{center}
\begin{tikzpicture}[scale=.6]
\draw[->] (-5.5,0) -- (5.5,0);
\draw[->] (0,-2.5) -- (0,2.5);
\foreach \x in {-5,...,-1,1,2,...,5}
\draw (\x,.1) -- (\x,-.1);
\foreach \y in {-2,-1,1,2}
\draw (.1,\y) -- (-.1,\y);
\draw (-5,-2) -- (5,2);
\end{tikzpicture}
\end{center}
\end{punkt}

\begin{punkt}
Likningen $xy = 1$ kan vi ikke skrive om til å bli på formen
$ax + by = c$, så den er ikke lineær.  Når vi skal tegne grafen, er
det igjen praktisk å skrive likningen slik at $y$ er alene:
\[
y = \frac{1}{x}
\]
Grafen blir slik:
\begin{center}
\begin{tikzpicture}[scale=1]
\begin{axis}[
axis y line=center,
axis x line=middle,
xmax=5,xmin=-5,
ymin=-5,ymax=5,
xtick={0},
ytick={0},
width=.5\textwidth,
anchor=center,
]
\addplot[black,domain=.2:5] {1/x} ;
\addplot[black,domain=-5:-.2] {1/x} ;
\end{axis}
\end{tikzpicture}
\end{center}
\end{punkt}

\begin{punkt}
Likningen
\[
\frac{x+y}{2} = \frac{7}{3}
\]
kan vi lett skrive om slik at vi ser at den er lineær:
\[
x + y = \frac{14}{3}
\]
Så kan vi sette $y$ alene --
\[
y = \frac{14}{3} - x
\]
-- og tegne grafen:
\begin{center}
\begin{tikzpicture}[scale=.6]
\draw[->] (-1.5,0) -- (5.5,0);
\draw[->] (0,-1.5) -- (0,5.5);
\foreach \x in {-1,1,2,...,5}
\draw (\x,.1) -- (\x,-.1);% node[anchor=north] {$\x$};
\foreach \y in {-1,1,2,...,5}
\draw (.1,\y) -- (-.1,\y);% node[anchor=east] {$\y$};
\draw (-.5,5.17) -- (5.5,-.83);
\end{tikzpicture}
\end{center}
\end{punkt}

\begin{punkt}
Likningen $x^2 = y^2$ ser veldig lite lineær ut.  Men så begynner du å
tenke på hvordan den kan forenkles, og da slår det deg antagelig at du
kan ta kvadratrot på begge sider og ende opp med bare
\[
x = y.
\]
Og det er jo så absolutt en lineær likning.

Men hvis du har tenkt slik, så må du stoppe opp et øyeblikk og tenke
deg om.  Både $x$ og~$y$ kan jo ha både positive og negative verdier,
men når vi opphøyer i andre, så blir resultatet alltid positivt.

Hvis $x^2$ for eksempel har verdien~$25$, så vet vi at $x$ må være
enten $5$ eller~$-5$, men vi kan ikke vite hvilken av dem det er.
Hvis vi setter $x = 5$ og $y = -5$, så er både $x^2$ og~$y^2$
lik~$25$, og dermed er likningen $x^2 = y^2$ oppfylt, selv om $x = y$
ikke er det.

Så det er ikke riktig å forenkle likningen $x^2 = y^2$ til $x = y$.
Men vi kan forenkle den til at vi har enten $x = y$ eller $x = -y$.
Da blir det lett å tegne grafen; den består av to rette linjer:
\begin{center}
\begin{tikzpicture}[scale=.4]
\draw[->] (-5,0) -- (5,0);
\draw[->] (0,-5) -- (0,5);
\draw (-5,-5) -- (5,5);
\draw (-5,5) -- (5,-5);
\end{tikzpicture}
\end{center}

Selv om grafen består av rette linjer, er ikke likningen lineær, siden
det ikke er mulig å skrive den på formen $ax + by = c$.  (Og hvis det
hadde vært mulig å skrive den på denne formen, så måtte grafen ha
blitt én rett linje, ikke to.)

(Vi kunne også ha valgt å skrive om likningen til $\abs{x} = \abs{y}$,
der notasjonen $\abs{x}$ står for absoluttverdien til~$x$.)
\end{punkt}

\begin{punkt}
Hva med likningen $x^2 + y^2 + 2 = 0$?  Det er ganske tydelig at den
ikke er lineær.  Men hvordan blir grafen?

Vi kan prøve å sette $y$ alene på den ene siden; da får vi:
\[
y = \pm \sqrt{-2 - x^2}
\]
Men dette ser rart ut.  Roten av minus~$2$ minus~$x^2$?  Hva skal det
bli?

Vel, det går ikke an.  Uttrykket $x^2$ blir jo alltid positivt
(eller~$0$), så $-2-x^2$ blir alltid negativt, og dermed kan vi ikke
ta kvadratroten av det.  (Så da kan du si at vi har \emph{rot}et oss
skikkelig bort.  He-he.)

La oss gå tilbake til den opprinnelige likningen.  Siden $x^2$
og~$y^2$ aldri kan bli negative, blir uttrykket på venstresiden aldri
mindre enn~$2$, og dermed kan det aldri bli lik~$0$.  Likningen har
altså ingen løsning.  Grafen blir dermed \ldots{} vel, ingenting:
\begin{center}
\begin{tikzpicture}[scale=.4]
\draw[->] (-5,0) -- (5,0);
\draw[->] (0,-5) -- (0,5);
\end{tikzpicture}
\end{center}
\end{punkt}
\end{losning}


\begin{oppgave}
Er disse to likningssystemene ekvivalente?
\[
\left\{
\begin{aligned}
  y &= x^2 - \tfrac{5}{4} \\
  x &= y + \tfrac{1}{2}
\end{aligned}
\right.
\qquad\qquad
\left\{
\begin{aligned}
  y &= \left( y + \tfrac{1}{2} \right)^2 - \tfrac{5}{4} \\
  y &= x^2 - \tfrac{5}{4}
\end{aligned}
\right.
\]
% y = y^2 + y + 1/4 - 5/4
% 0 = y^2 - 1
% y^2 = 1
% y = \pm 1
\end{oppgave}

\begin{losning}
I første omgang må vi vite nøyaktig hva som menes med at to
likningssystemer er ekvivalente.  På side~\pageref{def:ekvivalente}
ble dette definert slik:
\begin{quote}
Vi sier at to likningssystemer er \defterm{ekvivalente} dersom de har
samme løsninger.
\end{quote}
Det er altså dette vi må sjekke.

Det finnes minst to måter å finne ut dette på: Vi kan løse begge
systemene, og så sammenligne løsningene.  Eller vi kan sammenligne
likningene i de to systemene.  Dersom det er sånn at enhver likning i
det ene systemet er en logisk konsekvens av likningene i det andre, og
omvendt, så må systemene være ekvivalente.

I de aktuelle likningssystemene --
\[
\left\{
\begin{aligned}
  y &= x^2 - \tfrac{5}{4} \\
  x &= y + \tfrac{1}{2}
\end{aligned}
\right.
\qquad\qquad
\left\{
\begin{aligned}
  y &= \left( y + \tfrac{1}{2} \right)^2 - \tfrac{5}{4} \\
  y &= x^2 - \tfrac{5}{4}
\end{aligned}
\right.
\]
-- ser vi at vi får den første likningen i systemet til høyre ved å
kombinere de to likningene i systemet til venstre.  Og den andre
likningen i systemet til høyre er jo den samme som den første likningen
i systemet til venstre.

Vi har altså ganske lett funnet ut at systemet til høyre følger fra
systemet til venstre.  Da kan det være fristende å hoppe til
konklusjonen at de to systemene er ekvivalente.  Men hvis vi gjør det,
så har vi hoppet over noe viktig: Vi har ikke sjekket om systemet til
venstre følger fra systemet til høyre.

Og når vi prøver å finne ut av det, så blir det vanskeligere.  Det er
ikke åpenbart hvordan vi skal komme frem til likningen
\[
x = y + \frac{1}{2}
\]
ved å kombinere de to likningene i systemet til høyre.

Da kan vi prøve en annen strategi: Vi kan tegne grafene til likningene
i begge systemene, og sammenligne bildene vi får.

Her er grafene til likningene i det første systemet:
\begin{center}
\begin{tikzpicture}
\begin{axis}[
axis y line=center,
axis x line=middle,
xmax=5,xmin=-5,
ymin=-2,ymax=5,
xtick={0},
ytick={0},
width=.5\textwidth,
anchor=center,
]
\addplot[black,samples=31,domain=-2.3:2.3] {x^2-5/4)} ;
\addplot[black,domain=-4:4] {x-0.5} ;
\end{axis}
\end{tikzpicture}
\end{center}
Vi ser at grafene skjærer hverandre to steder, så det er altså to
løsninger av likningssystemet.  Vi ser på tegningen omtrent hvor de to
skjæringspunktene befinner seg, og med litt enkel regning kommer vi
frem til at de to løsningene er:
\[
\left\{
\begin{aligned}
x &= -\frac{1}{2} \\
y &= -1
\end{aligned}
\right.
\qquad\text{og}\qquad
\left\{
\begin{aligned}
x &= \frac{3}{2} \\
y &= 1
\end{aligned}
\right.
\]
%TODO sjekk at løsningene er riktige; ta med utregning?

Hva med det andre likningssystemet?  Likningen
\[
y = \left( y + \frac{1}{2} \right)^2 - \frac{5}{4}
\]
kan det være greit å forenkle litt før vi tegner grafen.  Ved å gange
ut parentesen får vi:
\[
y = y^2 + 2 \cdot \frac{1}{2} \cdot y + \frac{1}{4} - \frac{5}{4}
\]
Dette kan vi forenkle videre til
\[
0 = y^2 - 1,
\]
som er det samme som
\[
(y - 1)(y + 1) = 0.
\]
Denne likningen sier ganske enkelt at $y$ er enten $1$ eller~$-1$, og
$x$ kan være hva som helst.  Grafen dens består dermed av to rette
linjer:
\begin{center}
\begin{tikzpicture}
\begin{axis}[
axis y line=center,
axis x line=middle,
xmax=5,xmin=-5,
ymin=-2,ymax=2,
xtick={0},
ytick={0},
width=.5\textwidth,
anchor=center,
]
\addplot[black,domain=-4.8:4.8] {1} ;
\addplot[black,domain=-4.8:4.8] {-1} ;
\end{axis}
\end{tikzpicture}
\end{center}
Hvis vi tegner grafene til begge likningene i systemet sammen, så ser
det slik ut:
\begin{center}
\begin{tikzpicture}
\begin{axis}[
axis y line=center,
axis x line=middle,
xmax=5,xmin=-5,
ymin=-2,ymax=5,
xtick={0},
ytick={0},
width=.5\textwidth,
anchor=center,
]
\addplot[black,samples=31,domain=-2.3:2.3] {x^2-5/4)} ;
\addplot[black,domain=-4.8:4.8] {1} ;
\addplot[black,domain=-4.8:4.8] {-1} ;
\end{axis}
\end{tikzpicture}
\end{center}
Når vi skal se etter løsningene her, må vi passe på at vi tenker
riktig.  Det kan se ut som det ikke er noe sted der alle grafene
møtes, og dermed ingen løsning.  Men vi må huske på at de to rette
linjene til sammen utgjør grafen til den første likningen i systemet.

Det er fire punkter der de to grafene krysser hverandre, og dermed
fire løsninger av systemet.  Med litt utregning kan vi komme frem til
eksakte tall for løsningene:
% y = x^2 - 5/4 ; y = \pm 1
% x^2 = y + 5/4 = \pm 1 + 5/4
% x^2 = 9/4 eller x^2 = 1/4
% ( y = 1 ; x = \pm 3/2 ) eller ( y = -1 ; x = \pm 1/2 )
\[
\left\{
\begin{aligned}
x &= -\frac{1}{2} \\
y &= -1
\end{aligned}
\right.
\qquad
\left\{
\begin{aligned}
x &= \frac{1}{2} \\
y &= -1
\end{aligned}
\right.
\qquad
\left\{
\begin{aligned}
x &= -\frac{3}{2} \\
y &= 1
\end{aligned}
\right.
\qquad
\left\{
\begin{aligned}
x &= \frac{3}{2} \\
y &= 1
\end{aligned}
\right.
\]
Vi visste jo allerede at de to løsningene av det første
likningssystemet også måtte være løsninger av dette.  Men vi fikk
altså to løsninger til.  Det betyr at de to likningssystemene ikke har
de samme løsningene, og dermed er de ikke ekvivalente.
\end{losning}


\begin{oppgave}
Er disse to likningssystemene ekvivalente?
\[
\systeme{
  x + 2y - z = 0,
  2x + 2y = -3,
  4x + 8y - 3z = 2
}
\qquad
\systeme{
  3x + 4y - z = -3,
  x + y = -\frac{3}{2},
  5x + 10y - 4z = 2,
  7x + 12y - 4z = -1
}
\]
\end{oppgave}

\begin{losning}
Her må vi igjen enten løse begge likningssystemene (og sjekke om de
har samme løsninger) eller prøve å sammenligne likningene direkte.

Det kan se ut som det blir litt jobb å løse disse likningssystemene,
så latskapsimpulsen får oss til å heller prøve å sammenligne
likningene i de to systemene.  Da må vi bare prøve oss frem, og se om
vi kan utlede noen av likningene i det andre systemet ved å bruke de
som er i det første, og omvendt.

Etter å ha stirret litt på likningene, ser vi kanskje at det kan være
en god idé å prøve å legge sammen de to øverste likningene i det
første systemet:
\[
\begin{aligned}
\phantom{2}x + 2y - z &= 0 \\
2x + 2y \phantom{\hbox{}-z} &= -3 \\[4pt] \hline
\rule{0pt}{11pt}
3x + 4y - z &= -3
\end{aligned}
\]
Se der, ja!  Jammen klarte vi å utlede første likning i det andre
systemet.

Den andre likningen i det andre systemet er ganske grei å komme frem
til: Den er andre likning i det første systemet delt på to.

Så er det bare å prøve seg frem videre, og antagelig finner du etter
hvert ut at du kan legge sammen første og siste likning i det første
systemet, og at du kan legge sammen alle tre likningene.  Da får du de
to siste likningene i det andre systemet.

Dermed har vi funnet ut at alle likningene i det andre systemet følger
fra likningene i det første.  Klarer vi å gå motsatt vei også?  Prøv!

Dersom du får til å utlede alle likningene i det første systemet ved
hjelp av likningene i det andre, så har du vist at de to
likningssystemene er ekvivalente.
\end{losning}


\begin{oppgave}
Løs likningssystemet
\[
\systeme{
  4x + 4y + 16z = 0,
  3x + 4y + 12z = 5,
  2x + 2y + 5z = 1
}
\]
ved å bruke samme fremgangsmåte som i eksempel~\ref{ex:gausseliminasjon},
men uten å skrive opp hele likningssystemet på nytt hver gang --
skriv totalmatrisen for hvert system isteden.
\end{oppgave}

\begin{losning}
Totalmatrisen for systemet er:
\[
\begin{amatrix}{3}
4 & 4 & 16 & 0 \\
3 & 4 & 12 & 5 \\
2 & 2 & 5 & 1
\end{amatrix}
\]
Nå har vi noen valgmuligheter for hvordan vi vil gå frem, så det vi
kommer med her er bare vårt forslag.  Du kan ha gjort det på en annen
måte og fremdeles kommet frem til riktig løsning.

Det kan virke fornuftig å begynne med å gange den øverste likningen
med~$1/4$.  Da får vi denne totalmatrisen:
\[
\begin{amatrix}{3}
1 & 1 & 4 & 0 \\
3 & 4 & 12 & 5 \\
2 & 2 & 5 & 1
\end{amatrix}
\]
Så kan vi trekke fra $3$~ganger den første likningen i den andre, og
trekke fra $2$~ganger den første likningen i den tredje:
\[
\begin{amatrix}{3}
1 & 1 & 4 & 0 \\
0 & 1 & 0 & 5 \\
0 & 0 & -3 & 1
\end{amatrix}
\]
Herfra kan vi nokså enkelt få ut løsningen.  Men hvis vi vil, kan vi
også fortsette med å skrive om systemet for å gjøre det enda klarere
hva løsningen er.  Vi kan gange den nederste likningen med $-1/3$:
\[
\begin{amatrix}{3}
1 & 1 & 4 & 0 \\
0 & 1 & 0 & 5 \\
0 & 0 & 1 & -\frac{1}{3}
\end{amatrix}
\]
Og så kan vi trekke fra den andre likningen i den første:
\[
\begin{amatrix}{3}
1 & 0 & 4 & -5 \\
0 & 1 & 0 & 5 \\
0 & 0 & 1 & -\frac{1}{3}
\end{amatrix}
\]
Og til slutt trekke fra $4$~ganger den tredje likningen i den første:
\[
\begin{amatrix}{3}
1 & 0 & 0 & -\frac{11}{3} \\
0 & 1 & 0 & 5 \\
0 & 0 & 1 & -\frac{1}{3}
\end{amatrix}
\]
Da ser vi direkte at løsningen er:
\[
x = -\frac{11}{3}
\qquad
y = 5
\qquad
z = -\frac{1}{3}
\]
\end{losning}


% TODO finn lineært likningssystem med mer enn én løsning
% (med forslag til løsning: tre rette linjer)
